% Standalone fragment: input from the parent appendix; no section command.
% Completed frozen R84/R85 protocol; separate conditional hypotheses remain untested.
\subsection{Completed frozen R84/R85 protocol}
\label{app:completed-final-protocol}
The completed R84/R85 studies retained both predeclared configurations: TraceMemory (CORR25) was primary and BOTH25 was secondary. Their roles were fixed before independent generation. Given the native source and target fields evaluated at the same source-reference states and guidance, three scheduler nodes closest to $\sigma\in\{0.25,0.50,0.75\}$ provide channel-RMS prompt-response maps. Each map is clipped after its 10th--90th percentile rescaling; a degenerate map is set to zero. Their mean is $E$, and $P=1-0.75E$. The source-memory intervention uses all target image queries at unit gate unless the spatial-memory candidate explicitly pools $P$ over the packed $2\times2$ latent patch. No evaluation annotation is read by the method. The empirical mean of $P$ depends on the response distribution and is not fixed by the floor.

For the source-correction residual $c_i$, the primary local-release candidate applies $\lambda P\odot c_i$ without global RMS renormalization. The secondary spatial-memory configuration combines this correction with the pooled query gate. The energy-matched control applies $0.25\RMS(P\odot c_i)c_i/\RMS(c_i)$, with zero correction for a zero residual, matching the direct local rule's RMS at the same state and step; its trajectory and cumulative energy may subsequently differ. Fixed source K/V injection, active steps, reference schedule, guidance and prompts are shared where applicable. The three paired native probes cost six additional field evaluations per image; all cost values are recorded from actual calls.

\noindent\textbf{Development and independent evaluation.}
The development study contains ten previously used sources and 130 completed and automatically evaluated outputs: local-release coefficients $0.15$, $0.25$, and $0.35$, the combined spatial-memory candidate, the energy-matched control, and FlowEdit N28 each use two seeds; RF-Edit I08 uses seed 0. The completed independent study contains 40 new sources, four in each official category, selected by fixed hashing with source-content exclusion. Its six two-seed arms are local-release $0.25$, combined spatial memory, energy matching, uniform correction with source K/V, FlowEdit N26, and RF-Inversion ETA03; RF-Edit I08 uses seed 0. The final package extends the original 440-output independent plan by 80 ETA03 outputs while retaining every original job and source. The independent study therefore completed 520 automatically evaluated outputs, and the two studies total 650. ETA03 is an additional baseline comparison; the primary candidate and 40-source selection were fixed before generation. RF-Edit uses 25 steps and eight feature-injection steps. The previously studied 24 sources remain development data and are not pooled into the independent test.

For each image, the method receives a full-image update mask; the official edit mask is reserved for post-generation evaluation. The completed automatic evaluation reports target-prompt and directional CLIP as semantic proxies, preservation-region MAE, edit-region departure, DINO similarity, full-image and preservation LPIPS, NFE, and recorded runtime. The retained task-rating specification is given below. DINO and LPIPS characterize visual similarity or perceptual distance and do not establish requested-edit achievement. Before ratings, define a task-specific rubric and applicability flags from the target request. Blinded reviewers assess edit success, requested relations and materials, artifacts, and only those subject or pattern attributes the target requires to be preserved. Requested changes to category, pose, color, pattern or material are not preservation failures. An inapplicable flag is recorded as N/A, not as a success or failure, and is omitted from that task's conjunction: joint success requires edit success, no severe artifact, and every applicable preservation requirement to pass. Report applicable, N/A and rated counts by prespecified task category; retain all tasks, including the dog task, in the main table. A failed generation counts as unsuccessful in the planned joint denominator; uncompleted or unrated applicable items remain unknown. Aggregate seeds within source, compare matched source-level means, use source-bootstrap intervals and paired sign-flip tests with the preregistered Holm family. For deterministic RF-Edit, compare its single source result with the other arm's two-seed source mean; the two seeds are not independent sources. Empty preservation masks remain missing for that metric. An interval spanning zero is not an equivalence result.

\noindent\textbf{Conditional mechanism studies.}
If development results implicate support-dependent scaling, compare a small fixed-mask area or cap contrast with matched inputs and report actual local and global correction energies. The area-scaled coefficient $a\sqrt f\,\mathrm{RMS}(c)/\mathrm{RMS}(P c)$ equals the local coefficient $a$ only when $\mathrm{RMS}(Pc)=\sqrt f\,\mathrm{RMS}(c)$ and clipping is inactive; this need not hold for spatially heterogeneous residual energy or soft masks. If reproducible source-texture leakage motivates a source-pool study, run a small Full/Restricted source-pool $\times$ All/ROI target-query comparison with fixed pool construction, empty-pool handling, gate strength, seeds and cost. These conditional studies are separate from the 650-output completed final studies and do not establish outcomes in advance.

Cache reuse is valid only for identical state, timestep, prompt, guidance, schedule, model and attention-hook state. In particular, an inversion field under a different prompt or guidance is not interchangeable with a source-prompt probe, and a reference-state field is not generally the current edited-state field. Any lower-cost approximation requires matched quality and cost evaluation. The completed memory study provides native-pixel measurements for sixty new outputs. The local-release ablation and independent-source benchmark completed this frozen automatic-evaluation protocol; their outcomes appear in Section~\ref{sec:final-evaluation}. The conditional area-cap and source-pool studies remain separate hypotheses.


\noindent\textbf{Local attenuation and stability.}
For a fixed state and $0\leq P_j\leq1$, direct release obeys
\begin{equation}
\RMS(\lambda P\odot c_i)\leq\lambda\RMS(c_i),\qquad\lambda\geq0.
\end{equation}
Removing the global RMS restoration therefore attenuates the instantaneous correction relative to the corresponding full-image residual at that state. Different state trajectories can still exhibit different local energies and artifacts; per-step finite values, norms and decoded outcomes measure those effects. An area exponent and cap form an additional hypothesis, rather than a necessary consequence of removing RMS restoration.

The final automatic-evaluation pipeline used loaded CLIP, DINOv2 and LPIPS models. Preservation LPIPS averages the spatial LPIPS map over the independently specified preservation mask; its feature receptive fields may cross the mask boundary. This is reported separately from full-image LPIPS. Generation and CPU evaluation used separate stages. Automatic evaluation completed for all 650 outputs; empty preservation complements are excluded from regional denominators. The task-rating specification above is retained as the frozen rubric; its outcome coverage is stated in the main limitations.
